\documentclass[11pt]{article}
\usepackage[margin=1in]{geometry}
\usepackage[T1]{fontenc}
\usepackage{lmodern}
\usepackage{microtype}
\usepackage{amsmath,amssymb}

\begin{document}

\section{Experimental Setup}
\label{sec:experimental-setup}

\subsection{Factory layout}

Experiments use the three-floor factory generated by
\texttt{build\_three\_floors()} in \texttt{agv\_layouts.py}. It contains
three floor graphs with nine operational nodes per floor, for 27 operational
nodes total, connected by 38 edges. Four inter-floor edges represent lifts;
each lift has a configured length of 40~m. The fleet's single parking node is
$\pi=\texttt{1A}$. For configurations with lobbies, an additional lobby node
is inserted at each lift landing. Lobby nodes are waiting locations and are
excluded from the job-station sampling set. Thus, jobs are sampled from the
26 non-parking, non-lobby operational nodes.

\subsection{Jobs and fleet}

The generator \texttt{make\_jobs()} creates each job by sampling two distinct
stations uniformly from the 26 operational nodes: pickup $p_j$ and delivery
$p'_j$. Desired pickup time is drawn uniformly over the 600-second horizon,
\[
  t_j\sim U[0,600].
\]
For each job, an independent slack value is drawn as $s_j\sim U[30,200]$ and
the desired delivery time is set to
\[
  t'_j=t_j+d(p_j,p'_j)+\eta+\eta'+s_j,
\]
where $d(p_j,p'_j)$ is the shortest-path travel time on the factory graph.
The implementation rounds generated times to one decimal place. The random
seed determines the job set; the comparison uses seeds 1 through 10.

The main experiment uses 20 vehicles named \texttt{AGV-01} through
\texttt{AGV-20}, with fixed priority ranks $V_1$ through $V_{20}$ in that
order. The fleet-size sweep repeats the experiment with 4, 8, 12, 16, and 20
AGVs. Unless otherwise specified, the main comparison uses 100 jobs and 20
AGVs per seed.

\subsection{Planning parameters and configurations}

All configurations use loading time $\eta=16$~s, unloading time
$\eta'=16$~s, safety separation $\delta=10$~s, fitness weight $\lambda=1.0$,
and vehicle speed $v=1$~m/s. Lift links use their configured 40~m length, so
their non-stop travel time is 40~s at the specified speed.

The comparison script \texttt{agv\_compare.py} evaluates three configurations
on the same generated job sets:
\begin{description}
  \item[A --- Centralized paper.] The centralized prioritized planner with
        waiting permitted at the parking node only; lift landings have no
        exempt waiting lobbies.
  \item[B --- Centralized with lobbies.] The centralized planner uses the
        same route and conflict-planning procedure, with lift-lobby nodes
        added as exempt waiting locations.
  \item[C --- Hybrid.] The hybrid central-plus-floor-server planner uses the
        same lobbies. The central server assigns jobs and books lift trips;
        floor servers plan floor legs in first-come-first-served order
        (\texttt{order="fcfs"}).
\end{description}

The main comparison contains 10 seeds (1--10), 100 jobs, and 20 AGVs for each
of A, B, and C. A separate fleet sweep evaluates 4, 8, 12, 16, and 20 AGVs
over the same seeds and job-generation settings. For correctness checking of
the decomposition, the hybrid planner can also be run in priority order;
that check compares its plans with configuration B.

\subsection{Evaluation metrics}

The comparison records the following measures:
\begin{itemize}
  \item \textbf{Makespan:} time at which the last job delivery is completed.
  \item \textbf{Tardiness:} total positive delivery lateness across jobs,
        $\sum_j\max(0,C_j-t'_j)$, where $C_j$ is actual delivery completion.
  \item \textbf{Load imbalance:} maximum minus minimum AGV usage time, where
        usage time is measured until return to parking (zero for an AGV serving
        no jobs).
  \item \textbf{Waiting:} total and maximum per-AGV waiting time at parking
        or lift lobbies.
  \item \textbf{Compute time:} total planning CPU time across servers and CPU
        time at the busiest server; wall-clock planning time is also recorded.
  \item \textbf{Server state:} maximum number of reservation windows held by
        any one server.
  \item \textbf{Communication:} central-to-floor message count and the number
        of window records in the transmitted payload.
  \item \textbf{Conflict-resolution effort:} number of Algorithm~3 calls.
\end{itemize}

The comparison reports means over the ten seeds; the figures also show the
seed range where applicable. Fleet-sweep results report the corresponding
metrics by fleet size.

\subsection{Plan verification}

Every generated schedule is passed to the independent \texttt{verify()}
checker in \texttt{agv\_core.py}. The checker evaluates the completed plans
for start and end at $\pi$, node conflict separation (with parking/lobbies
exempt), bidirectional track conflict separation, no pause during edge
traversal, correct node-window dwell duration, and exactly-once job service.
For the main A/B/C comparison, a run is accepted only when the verifier
returns no violations. Configuration C in priority order is additionally
compared against B to check that the server decomposition preserves the
centralized-with-lobbies plans on the tested seeds.

\subsection{Software and execution environment}

The repository specifies Python 3.11 or 3.12 and the dependencies
\texttt{networkx}, \texttt{matplotlib}, \texttt{numpy}, and \texttt{pillow},
installed from \texttt{requirements.txt}. The requirements file names these
packages but does not pin their individual versions. The repository does not
record processor, memory, operating-system, or accelerator specifications;
therefore, hardware details and exact installed package versions are not
reported here. Runtime measurements should be interpreted in light of this
missing machine-level metadata.

\end{document}
